% Figure: projectile trajectories in a uniform field (no drag) at three
% launch angles, same launch speed, showing the 45-degree range maximum
% and the symmetry of complementary angles.
\begin{figure}[htbp]
\centering
\begin{tikzpicture}
\begin{axis}[
  width=12cm, height=7cm,
  xlabel={horizontal distance $x$ (\si{\meter})},
  ylabel={height $y$ (\si{\meter})},
  xmin=0, xmax=95, ymin=0, ymax=48,
  axis lines=left,
  samples=200,
  legend pos=north east,
  legend cell align=left,
  grid=major, grid style={enggray!20},
  tick label style={font=\scriptsize},
  label style={font=\small},
  legend style={font=\scriptsize},
]
% v0 = 30 m/s, g = 9.81. Trajectory y(x) = x tan(theta) - g x^2 / (2 v0^2 cos^2 theta).
% theta = 30: tan=0.5774, cos^2=0.75, g/(2 v0^2 cos^2) = 9.81/(2*900*0.75)=0.007267
\addplot[engblue, very thick, domain=0:79.4] {0.5774*x - 0.007267*x^2};
\addlegendentry{$\theta=30^{\circ}$}
% theta = 45: tan=1, cos^2=0.5, coeff = 9.81/(2*900*0.5)=0.0109
\addplot[engred, very thick, domain=0:91.7] {1.0*x - 0.0109*x^2};
\addlegendentry{$\theta=45^{\circ}$ (max range)}
% theta = 60: tan=1.7321, cos^2=0.25, coeff = 9.81/(2*900*0.25)=0.0218
\addplot[engamber, very thick, domain=0:79.4] {1.7321*x - 0.0218*x^2};
\addlegendentry{$\theta=60^{\circ}$}
\end{axis}
\end{tikzpicture}
\caption[Projectile trajectories at three launch angles]{Drag-free
projectile trajectories for a fixed launch speed $v_{0}=30\,\si{\meter
\per\second}$ at three launch angles. The range $R=v_{0}^{2}\sin
2\theta/g$ is maximised at $\theta=45^{\circ}$ (red), where it reaches
$R_{\max}=v_{0}^{2}/g\approx 91.7\,\si{\meter}$. The complementary
angles $30^{\circ}$ (blue) and $60^{\circ}$ (amber) share the same
range because $\sin 2\theta$ takes the same value at $\theta$ and
$90^{\circ}-\theta$; the steeper launch trades range for height and
flight time. Real projectiles peak well below $45^{\circ}$ once drag
is admitted, because the longer flight time of a high launch exposes
the body to drag for longer.}
\label{fig:vol03ch04:projectile-angles}
\end{figure}
